% =============================================================================
% Table 02 — Unity of Source and Aggregate
% Book I — Ontology of Viability
% Chapter 03 — Unity and Plurality
%
% Purpose:
%   To distinguish unity of source, unity of aggregate,
%   and unity of identity according to their respective
%   ontological subjects and relations.
%
% =============================================================================

\begin{table}[htbp]
\centering
\caption{Ontological Distinctions Between Unity of Source,
Unity of Aggregate, and Unity of Identity}
\label{tab:unity-source-aggregate-identity}

\renewcommand{\arraystretch}{1.45}
\setlength{\tabcolsep}{5pt}

\begin{tabular}{
    p{2.7cm}
    p{3.1cm}
    p{3.4cm}
    p{3.5cm}
}
\hline

\textbf{Form of Unity}
&
\textbf{Primary Subject}
&
\textbf{Ontological Basis}
&
\textbf{Limit of Inference}
\\

\hline

Unity of Source
&
Several manifestations or entities related to a common source
&
A relation of origin, grounding, or ontological derivation
&
Does not establish that all manifestations possess one identity
\\

\hline

Unity of Aggregate
&
Multiple components considered as an organized or bounded whole
&
Relations among components, including organization, integration, or composition
&
Does not establish that the aggregate necessarily constitutes one identity
\\

\hline

Unity of Identity
&
Identity considered in its ontological unity
&
The unity attributable to identity itself, requiring an account distinct from source and composition
&
Cannot be inferred solely from common origin, numerical singularity, structural integration, or invariant features
\\

\hline
\end{tabular}

\vspace{0.5em}

\begin{minipage}{0.94\textwidth}
\small
\textit{Conceptual qualification:}
The three forms of unity are distinguished according to
their respective ontological subjects and relations.
Their distinction does not imply that they are necessarily
independent in every existent, nor does it establish
a universal priority among them.
They may be related in particular ontological configurations,
but no one of them is sufficient, by itself, to establish
the unity of identity.
\end{minipage}

\end{table}

% =============================================================================
% End of Table 02
% =============================================================================